% ============================================================
\chapter{Applications}
\label{ch:applications}
% ============================================================

Geometric deep learning is not a theoretical exercise: it solves concrete problems that classical methods could not address. In chemistry, GNNs predict molecular properties with accuracy rivalling ab initio quantum mechanics, but a thousand times faster. In particle physics, they reconstruct trajectories in the detectors at CERN's LHC. In biology, AlphaFold (Jumper et al., 2021) uses geometric transformers to predict the 3D structure of proteins, solving a problem that had been open for fifty years. In social networks, GNNs detect communities and predict links. This chapter surveys these applications, showing how geometric principles translate into practical performance.

\begin{intuition}
Geometric deep learning finds applications in numerous domains where
data possesses a natural non-Euclidean structure: molecules (graphs),
social networks (graphs), surfaces (manifolds), physics simulations
(meshes and particles). This chapter presents the most impactful
applications.
\end{intuition}

% ============================================================
\section{Drug discovery}
\label{sec:drug_discovery}
% ============================================================

\begin{definition}[Molecular graph representation]
\index{molecular graph}
A molecule is represented as a graph $G = (V, E)$ where:
\begin{itemize}
  \item Nodes $v \in V$ represent atoms, with features $h_v$
        (atomic type, charge, hybridization, etc.).
  \item Edges $e \in E$ represent chemical bonds, with features
        $h_e$ (bond type, distance, etc.).
\end{itemize}
\end{definition}

\begin{example}[Molecular property prediction]
\index{molecular prediction}
The \textbf{MoleculeNet} benchmark (Wu et al.\ 2018) includes
several tasks:
\begin{enumerate}
  \item \textbf{ESOL}: aqueous solubility prediction (regression).
  \item \textbf{BBBP}: blood-brain barrier penetration (binary
        classification).
  \item \textbf{Tox21}: toxicity across 12 biological targets
        (multi-task classification).
\end{enumerate}
GNNs (SchNet, DimeNet, GemNet) achieve the best published results on these benchmarks thanks to their
ability to encode molecular structure.
\end{example}

\begin{definition}[Message passing for molecules]
For a molecular graph, SchNet's message passing
(Sch\"utt et al.\ 2017)\index{SchNet} integrates interatomic
distances:
\[
  m_{ij} = \phi_{\mathrm{msg}}(h_i, h_j, \norm{x_i - x_j})
\]
where $x_i \in \R^3$ are atomic positions. The continuous distance
filter makes the model invariant under rotation and translation.
\end{definition}

\begin{minted}{python}
import torch
from torch_geometric.nn import SchNet

# SchNet for molecular energy prediction
model = SchNet(
    hidden_channels=128,
    num_filters=128,
    num_interactions=6,
    num_gaussians=50,
    cutoff=10.0,
)

# z: atomic numbers, pos: 3D positions, batch: batch indices
# energy = model(z, pos, batch)
\end{minted}

% ============================================================
\section{Molecular property prediction}
\label{sec:molecular_properties}
% ============================================================

\begin{theorem}[Universality of equivariant GNNs for molecules]
\index{molecular universality}
An $\mathrm{E}(3)$-equivariant GNN with tensorial features of
sufficient degree can approximate any continuous equivariant
function on molecular graphs with atomic positions.
\end{theorem}

\begin{remark}
Recent models (DimeNet, SphereNet, PaiNN, MACE) use not only
distances but also angles and dihedrals to improve expressiveness.
\end{remark}

% ============================================================
\section{Social network analysis}
\label{sec:social_networks}
% ============================================================

\begin{definition}[Social network tasks]
\index{social network}
On a social network modeled as a graph:
\begin{enumerate}
  \item \textbf{Node classification}: predict the interests or
        community of a user.
  \item \textbf{Link prediction}: predict future connections
        between users.
  \item \textbf{Community detection}: cluster the nodes.
\end{enumerate}
\end{definition}

\begin{example}[Node classification on Cora]
The \textbf{Cora}\index{Cora} dataset contains 2,708 scientific
papers (nodes) connected by citations (edges). Each paper must be
classified among 7 categories. A 2-layer GCN achieves $\sim 81\%$
accuracy with only $20$ labels per class.
\end{example}

\begin{proposition}[Over-smoothing]
\index{over-smoothing}
For a GCN with $L$ layers, node embeddings converge to a subspace
of dimension at most $1$ as $L \to \infty$:
\[
  \lim_{L \to \infty} h_v^{(L)} = c \cdot \sqrt{d_v} \quad
  \text{for all } v \in V
\]
where $d_v$ is the node degree. This is the \textbf{over-smoothing}
phenomenon that limits GNN depth.
\end{proposition}

% ============================================================
\section{Recommender systems}
\label{sec:recommender}
% ============================================================

\begin{definition}[Collaborative filtering on bipartite graphs]
\index{recommender system}
A recommender system can be modeled as a bipartite graph
$G = (U \cup I, E)$ where $U$ are users, $I$ are items, and
$(u, i) \in E$ if user $u$ interacted with item $i$.
\end{definition}

\begin{example}[PinSage --- Ying et al.\ 2018]
\index{PinSage}
\textbf{PinSage} is a GNN deployed by Pinterest on a graph of
3 billion nodes and 18 billion edges. It uses:
\begin{itemize}
  \item \textbf{Neighbor sampling}: importance-weighted random walks.
  \item \textbf{Local aggregation}: similar to GraphSAGE.
  \item \textbf{Mini-batch training}: on sampled subgraphs.
\end{itemize}
\end{example}

% ============================================================
\section{Physics simulation}
\label{sec:physics}
% ============================================================

\begin{definition}[Learned simulator via GNN]
\index{physics simulation}
A learned simulator models a physical system as an interaction
graph:
\begin{itemize}
  \item Nodes represent particles (position, velocity, type).
  \item Edges represent interactions (nearest neighbors).
\end{itemize}
The GNN predicts the acceleration of each particle:
\[
  \ddot{x}_i = f_\theta\!\left(h_i,
  \{h_j, x_j - x_i\}_{j \in \mathcal{N}(i)}\right).
\]
\end{definition}

\begin{example}[GNS --- Sanchez-Gonzalez et al.\ 2020]
\index{GNS}
The \textbf{Graph Network Simulator} (GNS) simulates fluids,
granular materials, and deformable solids. It generalizes to
unseen initial conditions and different domain geometries.
\end{example}

\begin{theorem}[Energy conservation by construction]
\index{Hamiltonian GNN}
By parameterizing a GNN as a Hamiltonian $H_\theta$ and integrating
Hamilton's equations:
\[
  \dot{q}_i = \frac{\partial H_\theta}{\partial p_i}, \qquad
  \dot{p}_i = -\frac{\partial H_\theta}{\partial q_i}
\]
one obtains a simulator that conserves energy by construction
(Hamiltonian GNN).
\end{theorem}

% ============================================================
\section{Weather prediction}
\label{sec:weather}
% ============================================================

\begin{example}[GraphCast --- Lam et al.\ 2023]
\index{GraphCast}
\textbf{GraphCast} (DeepMind) uses a GNN on an icosahedral mesh
of Earth for 10-day weather forecasting. It outperforms ECMWF's
HRES numerical model on $90\%$ of metrics, with computation time
of \textbf{less than one minute} versus hours for physical models.

The architecture uses:
\begin{enumerate}
  \item An \textbf{encoder} from grid to multi-resolution mesh.
  \item A \textbf{processor} GNN on the mesh (16 message-passing
        layers).
  \item A \textbf{decoder} from mesh back to grid.
\end{enumerate}
\end{example}

\begin{remark}
The success of GraphCast illustrates the potential of GNNs for
problems defined on non-planar domains (Earth's sphere).
\end{remark}

% ============================================================
\section{Exercises}
% ============================================================

\begin{exercise}
For a benzene molecule ($C_6H_6$), construct the molecular graph
(nodes, edges, features). How many messages are exchanged in one
iteration of message passing?
\end{exercise}

\begin{exercise}
Explain why permutation invariance of nodes is essential for
molecular graphs. What would happen if atoms were numbered
differently?
\end{exercise}

\begin{exercise}[Over-smoothing]
For a complete graph $K_n$ with a GCN without bias or nonlinearity,
show that features converge to a constant vector after $L$ layers
when $L$ is large.
\end{exercise}

\begin{exercise}[Implementation]
Implement a GNN for solubility prediction on the ESOL dataset
with PyTorch Geometric. Compare GCN, GAT, and SchNet.
\end{exercise}

\begin{exercise}[Simulation]
Implement a simple GNS to simulate $N$ particles interacting via a
Lennard-Jones potential. Compare with classical Verlet integration.
\end{exercise}
